\documentclass[11pt]{article}
\usepackage{amsmath,amssymb,amsthm}
\newtheorem{theorem}{Theorem}
\title{A Metric Audit of Finite Cauchy Token Features}
\author{}
\date{November 24, 2026}
\begin{document}
\maketitle
\paragraph{Scope.}
This written lemma concerns one proposed random-fingerprint argument.
It is not a head-complexity lower bound, and it does not rule out
small-margin or interacting-token constructions.

\begin{theorem}
Fix distinct real numbers $D_1,\ldots,D_L$ and an interval $I$
containing none of them. Let $T:\mathbb R^L\to\mathbb R^d$ be any
linear map, and put
\[
 v(t)=T\big((D_1-t)^{-1},\ldots,(D_L-t)^{-1}\big)^{\mathsf T}.
\]
For ordered points $t_1<\cdots<t_K$ in $I$ such that $v(t_j)\ne0$,
let $u_j=v(t_j)/\|v(t_j)\|_2$. Then
\[
 \sum_{j=1}^{K-1}\arccos\langle u_j,u_{j+1}\rangle
                  \le \pi(L-1). \tag{1}
\]
Consequently, if these unit vectors are pairwise separated by angle
at least $\theta>0$, then
\[
                    K\le 1+\pi(L-1)/\theta. \tag{2}
\]
The same feature columns multiplied by arbitrary nonzero real scalars
have at most $2+2\pi(L-1)/\theta$ pairwise $\theta$-separated normalized
directions.
\end{theorem}
\begin{proof}
Take a standard Gaussian vector $g\in\mathbb R^d$. For each specified
nonzero $v(t_j)$, the event $g\cdot v(t_j)=0$ has probability zero.
Outside the union of these finitely many null events, every sampled
value is nonzero. Also,
\[
 g\cdot v(t)=\sum_{l=1}^L\frac{a_l}{D_l-t}
   =\frac{P_g(t)}{\prod_{l=1}^L(D_l-t)},\qquad \deg P_g\le L-1.
\]
The denominator has constant nonzero sign on $I$. The polynomial is
not identically zero, because at least one of its specified values is
nonzero. A sign change between consecutive sampled values forces a
distinct root of $P_g$ in the corresponding open interval. Thus the
number of such sign changes is at most $L-1$.

For any two nonzero vectors $v,w$, rotational symmetry of a standard
Gaussian gives
\[
 \mathbb P\{\operatorname{sgn}(g\cdot v)\ne
                  \operatorname{sgn}(g\cdot w)\}
       =\frac{\arccos\langle v/\|v\|,w/\|w\|\rangle}{\pi}.
\]
Indeed, if the vectors are linearly independent, project onto their
plane. The angle of the projection is uniform, and the two opposite
wedges on which the signs differ have total angle twice the angle
between the vectors. Parallel and antiparallel vectors give the same
formula directly. Summing these probabilities over consecutive
samples and using the preceding deterministic root bound proves (1).
Every consecutive angle in a pairwise separated collection is at
least $\theta$, which gives (2). Finally, divide arbitrary scalar
multipliers into their positive and negative classes. In each class,
normalization either preserves all directions or reverses all of
them, so (2) applies separately and the two bounds add.
\end{proof}

\paragraph{Consequence for the proposed mechanism.}
Suppose a token $t$ is represented by its values at $L$ common exterior
Cauchy probes, followed by an arbitrary fixed linear transformation
and a token-dependent scalar residue. A family of $H$ tokens cannot
have a fixed positive pairwise angular separation when $L=O(\log H)$.
It needs $L=\Omega(H)$. Thus a random-hyperplane fingerprint proof
cannot obtain $H\log H$ independent bits merely by claiming that these
$L$-probe token directions behave like an exponentially large
spherical code. The argument leaves open fingerprints based on
vanishing separation, query-dependent cancellations, or genuinely
joint token invariants.

\paragraph{Verification boundary.}
The proof is written only. It changes no model definition, adds no
Lean declaration, and proves no asymptotic bound for worst-case head
complexity.
\end{document}
